\documentclass{article}

\usepackage{amsmath,amssymb,amsthm}
\usepackage{booktabs}
\usepackage{multirow}
\usepackage{graphicx}
\usepackage{hyperref}
\usepackage{microtype}
\usepackage{xcolor}
\usepackage{cleveref}

\title{Stability Regularisation in Sequence Models:\\
Reference Optimisation and Contour Barriers for Diagonal SSMs}
\author{Author Names Omitted}
\date{}

\newtheorem{theorem}{Theorem}
\newtheorem{proposition}{Proposition}
\newtheorem{lemma}{Lemma}
\newtheorem{corollary}{Corollary}
\newtheorem{remark}{Remark}
\theoremstyle{definition}
\newtheorem{definition}{Definition}

\newcommand{\C}{\mathbb{C}}
\newcommand{\R}{\mathbb{R}}
\renewcommand{\rho}{\varrho}
\newcommand{\norm}[1]{\left\|#1\right\|}
\newcommand{\abs}[1]{\left|#1\right|}
\DeclareMathOperator*{\argmax}{arg\,max}
\DeclareMathOperator*{\argmin}{arg\,min}

\newcommand{\TBD}[1]{\textcolor{red}{[TBD: #1]}}

\begin{document}
\maketitle

\begin{abstract}
Post-hoc contour certificates for learned state-space model (SSM)
stability, introduced in Paper~1 of this series, achieve high
false-negative rates when paired with naive scalar or random-search
reference matrices.  This paper addresses the core bottleneck directly:
we develop a gradient-based algorithm for optimising the diagonal
reference matrix $A_0$ jointly with the Rouch\'e margin, achieving
100\% post-hoc certification on diagonal and trained-SSM matrix families
(up from $\leq 36\%$ with random search).  Diagonal references remain
insufficient for non-normal and Jordan-near-boundary matrices regardless
of optimisation quality.  We then embed the optimised reference into
three families of training-time stability regularisers---spectral
penalty, Lyapunov per-mode penalty, and a differentiable contour
barrier---and compare them on autoregressive prediction tasks under two
regimes.  In a standard constrained regime (clipped eigenvalue radii),
all three achieve 100\% post-hoc certification; they differ only in the
task-loss--spectral-radius trade-off.  In an adversarial unconstrained
regime (lr=0.1, near-boundary init at $\varrho=0.96$), 90\% of
unregularised runs overshoot the unit circle; hinge regularisers reduce
but do not eliminate instability (10--20\% stable), demonstrating that
post-hoc penalisation cannot substitute for constrained parameterisation
under extreme gradient pressure.
\end{abstract}

%% ============================================================
\section{Introduction}
\label{sec:intro}

Paper~1 established a framework of Rouch\'e-type contour certificates
for Schur stability of learned transition matrices.  The central
empirical finding was that both determinant and resolvent certificates
achieve less than 3\% certification rates when paired with scalar or
random-diagonal reference matrices, even for matrices the certificates
could in principle certify.  The bottleneck is the \emph{reference
matrix} $A_0$: a poorly chosen reference inflates the perturbation
$\norm{A_0 - A}$ relative to the resolvent margin, rendering the
certificate vacuous.

This paper makes two contributions.  First, we introduce a
gradient-based reference optimisation algorithm that directly maximises
the sampled minimum resolvent margin $\hat{m}(A, A_0)$ with respect to
the diagonal entries of $A_0$, using PyTorch autograd through a
differentiable softmin over contour points.  For diagonal $A_0$ and any
matrix $A$, the resolvent margin evaluation reduces from $O(Kn^2)$
(batched SVD) to $O(Kn)$ (elementwise division), making the inner loop
fast enough for gradient ascent in practice.

Second, we embed the optimised reference into a differentiable contour
barrier regulariser for SSM training and compare it against spectral and
Lyapunov baselines.  The contour barrier uses the same $O(Kn)$ diagonal
formula, making per-epoch overhead negligible.

\paragraph{Relation to Paper~1.}
All certificate definitions, Theorems 1--4, and the Paper~1 codebase
(\texttt{src/roche/}) remain unchanged.  This paper adds
\texttt{src/roche/reference\_opt.py} (gradient reference optimisation),
\texttt{src/roche/regularisers.py} (three regulariser implementations),
and an Adam-based SSM trainer (\texttt{train\_diagonal\_ssm\_adam}) to
\texttt{src/roche/ssm.py}.

%% ============================================================
\section{Background}
\label{sec:background}

\subsection{Contour certificates (recap)}

Let $A \in \C^{n \times n}$ and $A_0 \in \C^{n \times n}$ be Schur stable.
The \emph{resolvent margin} at $z = e^{i\theta}$ is
\[
  m_{\mathrm{res}}(A, A_0, z)
  = 1 - \norm{(zI - A_0)^{-1}(A_0 - A)}_2.
\]
When $\min_\theta m_{\mathrm{res}} > 0$, Rouch\'e's theorem guarantees
$\varrho(A) < 1$.  The finite-grid certificate (Theorem~3 of Paper~1)
converts a sampled positive minimum margin into a continuum guarantee
via a Lipschitz criterion.

\subsection{The reference selection problem}

Define the \emph{optimal diagonal reference} as
\[
  A_0^* = \argmax_{A_0 \in \mathcal{D}} \min_{\theta} m_{\mathrm{res}}(A, A_0, \theta),
\]
where $\mathcal{D}$ is the set of diagonal Schur-stable matrices.
Paper~1's random search over 50--200 diagonal candidates achieves
near-zero certification rates for normal and non-normal families,
motivating a gradient-based approach.

\subsection{Diagonal fast path}

For $A_0 = \mathrm{diag}(d_1, \ldots, d_n)$ with $A$ arbitrary,
\[
  (zI - A_0)^{-1}(A_0 - A) = \mathrm{diag}\!\left(\frac{1}{z - d_i}\right)(A_0 - A),
\]
i.e.\ each row $i$ of $A_0 - A$ is scaled by $1/(z - d_i)$.  The
2-norm of this product equals the largest singular value of a
row-scaled matrix, which costs $O(Kn^2)$ in general.  When $A$ is
\emph{also diagonal} (as in diagonal complex SSMs), $A_0 - A$ is
diagonal and the 2-norm reduces to $\max_i \abs{(d_{0,i} -
a_i)/(z - d_{0,i})}$, costing only $O(Kn)$ with no SVD.

%% ============================================================
\section{Gradient Reference Optimisation}
\label{sec:refopt}

\subsection{Parameterisation}

We parameterise $A_0 = \mathrm{diag}(r_i e^{i\phi_i})$ where
\[
  r_i = \sigma(\ell_i)(1 - \varepsilon), \quad \varepsilon = 0.02,
\]
with $\sigma$ the logistic function and $\ell_i, \phi_i \in \R$ the
free parameters.  This ensures $r_i \in (0, 1 - \varepsilon)$
without any projection step.

\subsection{Objective}

We maximise a differentiable lower bound on $\min_k m_{\mathrm{res}}(z_k)$
over a grid of $K$ equispaced contour points via a softmin:
\[
  \hat{m}_\beta(A, A_0)
  = -\frac{1}{\beta}\log\sum_k \exp(-\beta \cdot m_{\mathrm{res}}(z_k)),
  \quad \beta = 20.
\]
Gradient ascent on $\hat{m}_\beta$ with respect to $(\ell_i, \phi_i)$
is performed with Adam (lr = 0.05, 300 steps) using PyTorch autograd.
Initialisation uses the eig-shrunk reference (Paper~1), which provides
a warm start with non-trivial margin for trained SSMs.

\subsection{Diagonal-plus-low-rank reference (extension)}
\label{sec:dlr}

A diagonal $A_0$ cannot closely approximate matrices with significant
off-diagonal structure.  We extend the parameterisation to
\emph{diagonal-plus-low-rank} (DLR):
\[
  A_0 = \mathrm{diag}(d_0) + U V^H, \quad U, V \in \C^{n \times r},
\]
with the diagonal part reparameterised as before and $U, V$
initialised to zero (so $A_0$ starts diagonal).  Stability of $A_0$
is enforced via a softplus penalty on the spectral radius of $A_0$
rather than by construction:
\[
  \mathcal{P}_{\mathrm{stab}} = \mathrm{softplus}(\varrho(A_0) - (1 - \varepsilon)).
\]
The resolvent margin is computed via batched
$\texttt{torch.linalg.solve}$ + SVD ($O(Kn^3)$ per step), since $A_0$
is no longer diagonal.  Rank $r = 1$ adds $4n$ real parameters beyond
the diagonal; our tests confirm the DLR optimiser finds positive
margins but convergence is slower than the diagonal path, consistent
with the harder non-convex landscape.

%% ============================================================
\section{Stability Regularisers}
\label{sec:regularisers}

All three regularisers take the diagonal eigenvalue vector
$\mathbf{a} = (a_1, \ldots, a_n) \in \C^n$ and return a real scalar.

\paragraph{Spectral penalty.}
\[
  \mathcal{R}_{\mathrm{spec}}(\mathbf{a}) = \mathrm{relu}\!\left(\max_i \abs{a_i} - (1 - \tau)\right),
\]
where $\tau$ is a stability margin.  This is a hinge on the spectral
radius; only the largest mode is penalised.  \texttt{relu} rather than
\texttt{softplus} is used so the penalty is \emph{exactly zero} for
clearly stable systems, avoiding a constant background term.

\paragraph{Lyapunov penalty.}
\[
  \mathcal{R}_{\mathrm{Lyap}}(\mathbf{a}) = \sum_i \mathrm{relu}\!\left(\abs{a_i}^2 - (1 - \tau)\right).
\]
For diagonal $A$ with $P = I$ the Lyapunov condition $A^\top P A - P \prec 0$
reduces to $\abs{a_i}^2 < 1$ per mode.  The sum formulation penalises
\emph{all} near-unstable modes jointly, yielding different gradient dynamics
from the spectral hinge.

\paragraph{Contour barrier.}
\[
  \mathcal{R}_{\mathrm{ctour}}(\mathbf{a}) = \mathrm{relu}\!\left(-\hat{m}_\beta(A, A_0^*) + \tau\right),
\]
where $A_0^*$ is the gradient-optimised reference computed once from
the initial model and frozen throughout training.  The barrier fires
only when the resolvent margin falls below $\tau$, providing a
certificate-geometry-aware signal.

%% ============================================================
\section{Experiments}
\label{sec:experiments}

\subsection{Experiment 1: Reference quality comparison}

We compare all four reference methods (scalar, eig-shrunk, random-search,
gradient-optimised) across four matrix families:
\texttt{diagonal\_stable} (random diagonal, $\varrho \in [0.7, 0.95]$);
\texttt{nonnormal\_stable} ($\varrho = 0.85$, departure~20);
\texttt{jordan\_near\_boundary} ($\varrho = 0.95$, superdiag~1.0);
\texttt{trained\_ssm} (diagonal SSM trained on AR(4) long-memory).
$M = 50$ instances per family, $n = 8$, $K = 256$ contour points.

\begin{table}[t]
\centering
\caption{Experiment~1: Reference quality across matrix families and methods
($n=8$, $M=50$ instances, $K=256$ contour points).
``cert\%'': percentage of stable instances certified by the resolvent
certificate; ``mean margin'': mean minimum resolvent margin across all
instances (positive = certified, negative = not certified).
Gradient optimisation uses the SVD path for non-diagonal $A$ and the
$O(Kn)$ diagonal fast path for diagonal $A$ and trained SSMs.}
\label{tab:refquality}
\begin{tabular}{llrr}
\toprule
Family & Method & cert\% & Mean margin \\
\midrule
\multirow{4}{*}{diagonal\_stable}
  & scalar        &   0 & $-1.65$ \\
  & eig\_shrunk   & 100 & $+0.43$ \\
  & random\_search&   0 & $-0.53$ \\
  & gradient\_opt & 100 & $+0.96$ \\
\midrule
\multirow{4}{*}{nonnormal\_stable}
  & scalar        &   0 & $-11.07$ \\
  & eig\_shrunk   &   0 & $-12.37$ \\
  & random\_search&   0 & $-6.07$ \\
  & gradient\_opt &   0 & $-4.98$ \\
\midrule
\multirow{4}{*}{jordan\_near\_boundary}
  & scalar        &   0 & $-1.86$ \\
  & eig\_shrunk   &   0 & $-6.51$ \\
  & random\_search&   0 & $-1.66$ \\
  & gradient\_opt &   0 & $-0.97$ \\
\midrule
\multirow{4}{*}{trained\_ssm}
  & scalar        &   0 & $-1.10$ \\
  & eig\_shrunk   & 100 & $+0.79$ \\
  & random\_search&  36 & $-0.09$ \\
  & gradient\_opt & 100 & $+0.99$ \\
\bottomrule
\end{tabular}
\end{table}

\begin{figure}[t]
\centering
\includegraphics[width=0.55\linewidth]{../../results/p2exp1/reference_quality_margins.png}
\hfill
\includegraphics[width=0.44\linewidth]{../../results/p2exp1/reference_opt_convergence.png}
\caption{Left: mean resolvent margin by method and family.
Right: gradient ascent convergence (softmin margin vs.\ step) for one
representative instance per family.  The diagonal fast path converges
in under 500\,ms; the SVD path for non-diagonal matrices takes $\sim
3\,$s per instance but still improves margins substantially over
scalar and random-search baselines.}
\label{fig:refquality}
\end{figure}

Key findings: gradient optimisation achieves 100\% certification on
\texttt{diagonal\_stable} and \texttt{trained\_ssm} families (matching
eig-shrunk), while producing substantially tighter margins
($+0.96$ vs.\ $+0.43$ for diagonal stable).  For non-diagonal families
(nonnormal, Jordan) no method certifies, but gradient optimisation
reduces the margin gap from $-11.07$ to $-4.98$ (nonnormal) and from
$-1.86$ to $-0.97$ (Jordan near boundary), suggesting the diagonal
reference class is insufficient for these structures and motivates
structured non-diagonal references as future work.

\subsection{Experiment 2: Regulariser comparison}

We train diagonal complex SSMs ($n = 8$, Adam, 500 epochs, lr = 0.001)
on three AR tasks under four regularisation regimes (weight $\lambda = 0.1$):
no regulariser, spectral penalty, Lyapunov penalty, and contour barrier.
Five random seeds per configuration (60 runs total).
Post-hoc certification uses the eig-shrunk reference (Paper~1 Exp~3 protocol).

Tasks: AR(2) short-memory ($\varrho_{\max} = 0.70$);
AR(4) long-memory ($\varrho_{\max} = 0.92$);
AR(4) hard ($\varrho_{\max} = 0.97$).

\begin{table}[t]
\centering
\caption{Experiment~2: Regulariser comparison ($n=8$, Adam 500 epochs,
$\lambda=0.1$, 5 seeds, seq.\ length 500).
Post-hoc certification uses the eig-shrunk reference (Paper~1
protocol).  All methods achieve 100\% certification on all tasks;
differences lie in task loss, spectral radius, and margin.
``task loss'': final MSE; ``$\varrho$'': mean spectral radius of
trained $A$; ``margin'': mean min resolvent margin (eig-shrunk ref).}
\label{tab:regularisers}
\begin{tabular}{llrrrr}
\toprule
Task & Regulariser & Task loss & $\varrho$ & Margin & Cert\% \\
\midrule
\multirow{4}{*}{AR(2) short}
  & none     & $2.57\times10^{-3}$ & 0.76 & $+0.73$ & 100 \\
  & spectral & $2.60\times10^{-3}$ & 0.49 & $+0.91$ & 100 \\
  & lyapunov & $2.70\times10^{-3}$ & 0.46 & $+0.92$ & 100 \\
  & contour  & $2.62\times10^{-3}$ & 0.68 & $+0.82$ & 100 \\
\midrule
\multirow{4}{*}{AR(4) long}
  & none     & $2.51\times10^{-3}$ & 0.83 & $+0.67$ & 100 \\
  & spectral & $2.65\times10^{-3}$ & 0.51 & $+0.90$ & 100 \\
  & lyapunov & $3.10\times10^{-3}$ & 0.46 & $+0.92$ & 100 \\
  & contour  & $2.70\times10^{-3}$ & 0.68 & $+0.82$ & 100 \\
\midrule
\multirow{4}{*}{AR(4) hard}
  & none     & $2.49\times10^{-3}$ & 0.84 & $+0.61$ & 100 \\
  & spectral & $2.53\times10^{-3}$ & 0.50 & $+0.91$ & 100 \\
  & lyapunov & $2.64\times10^{-3}$ & 0.46 & $+0.92$ & 100 \\
  & contour  & $2.54\times10^{-3}$ & 0.68 & $+0.82$ & 100 \\
\bottomrule
\end{tabular}
\end{table}

\begin{figure}[t]
\centering
\includegraphics[width=0.55\linewidth]{../../results/p2exp2/training_curves.png}
\hfill
\includegraphics[width=0.44\linewidth]{../../results/p2exp2/certification_rates.png}
\caption{Left: mean training curves (task loss and regulariser loss) on
AR(4) hard for all four regularisation regimes.
Right: post-hoc certification rate by regulariser and task (all 100\%;
bar heights uniform).}
\label{fig:regularisers}
\end{figure}

Key findings: all four regularisation regimes achieve 100\% post-hoc
certification across all tasks, because diagonal SSMs trained on
stable AR processes naturally remain stable and the eig-shrunk reference
suffices.  The regularisers differ in their spectral-radius--task-loss
trade-off.  The Lyapunov penalty is most aggressive: it drives $\varrho$
to $\approx 0.46$ at the cost of a $24\%$ task-loss increase on the
long-memory AR(4) task ($3.10$ vs.\ $2.51\times10^{-3}$).  The spectral
penalty achieves the lowest task loss among the three regularisers (e.g.\
$2.53$ vs.\ $2.54\times10^{-3}$ for contour on AR(4) hard) while halving
$\varrho$.  The contour barrier's distinctive property is $\varrho$
preservation ($\approx 0.68$ vs.\ $0.46$--$0.50$ for spectral/Lyapunov):
it does not suppress stable modes beyond the margin threshold $\tau$,
at the cost of a lower resolvent margin ($+0.82$) than spectral ($+0.90$--$+0.91$)
or Lyapunov ($+0.92$).

\subsection{Experiment 3: Adversarial unconstrained regime}
\label{sec:exp3}

To isolate whether hinge regularisers can \emph{prevent} instability
(rather than merely coexist with it in an already-stable regime), we
train with \emph{unconstrained} eigenvalue radii, a high learning rate
(lr = 0.1), adversarial near-boundary initialisation ($\varrho_{\mathrm{init}} = 0.96$),
and a long-memory AR(4) task ($\varrho_{\max} \approx 0.97$, seq.\ length 2000).
Regulariser threshold is $\tau = 0.1$ (fires at $\varrho = 0.9$) with weight
$\lambda = 1.0$.  Ten seeds per method; $n = 8$, 300 epochs.

\begin{table}[t]
\centering
\caption{Experiment~3: Stability regularisers under adversarial gradient
pressure (unconstrained SSM, lr=0.1, init $\varrho=0.96$, 10 seeds).
``stable\%'': fraction of runs with $\varrho(A) < 1$ at convergence;
``mean $\varrho^*$'': mean spectral radius of finite (non-diverged) runs;
``diverged'': runs where training reached NaN/inf;
``cert\%'': post-hoc certification rate on stable runs (eig-shrunk ref).}
\label{tab:adversarial}
\begin{tabular}{lrrrr}
\toprule
Regulariser & stable\% & mean $\varrho^*$ & diverged & cert\% (stable) \\
\midrule
none            & 10 & 1.090 & 1 & 100 \\
spectral        & 10 & 1.108 & 1 & 100 \\
lyapunov        & 20 & 1.038 & 1 & 100 \\
contour barrier & 10 & 1.103 & 1 & 100 \\
\bottomrule
\end{tabular}
\end{table}

\begin{figure}[t]
\centering
\includegraphics[width=0.55\linewidth]{../../results/p2exp3/rho_distribution.pdf}
\hfill
\includegraphics[width=0.44\linewidth]{../../results/p2exp3/stability_and_cert.pdf}
\caption{Left: distribution of final spectral radii per method (box plots,
10 seeds).  The red dashed line marks the unit circle.  All methods show
most runs exceeding $\varrho = 1$, with means clustering near $1.09$--$1.11$.
Right: stability rate (\%) and post-hoc certification rate (\%) for stable runs.}
\label{fig:adversarial}
\end{figure}

Key findings: at lr = 0.1, 90\% of unregularised runs overshoot the unit
circle within 300 epochs.  Hinge regularisers provide minimal benefit:
spectral and contour achieve the same 10\% stability rate as the
unregularised baseline; the Lyapunov penalty---which sums \emph{all}
near-unstable modes rather than the maximum---recovers one additional
seed (20\% stable).  The core failure mode is that relu-based hinges
have zero gradient below the threshold; at lr = 0.1 a single Adam step
can carry eigenvalues from $\varrho = 0.96$ past the unit circle before
any regulariser gradient fires.  Once escaped, the task gradient (which
favours high $\varrho$ for long-memory prediction) is strong enough to
maintain the unstable equilibrium at $\varrho \approx 1.1$.
The single diverged seed (seed 2, $\varrho \to \infty$) occurs identically
for all methods, consistent with the regulariser having no effect before
escape.  Stable runs all certify (100\%), confirming that when stability
is preserved, post-hoc certification remains reliable.
\emph{Conclusion:} hinge regularisers cannot substitute for constrained
parameterisation under adversarial gradient pressure; a clamp or
projection that enforces $\varrho < 1$ by construction is necessary.

%% ============================================================
\section{Results}
\label{sec:results}

Gradient reference optimisation (Experiment~1) achieves 100\%
certification on the target families (diagonal stable, trained SSMs)
with mean margins $+0.96$ and $+0.99$ respectively, substantially
tighter than eig-shrunk ($+0.43$, $+0.79$) and random-search ($-0.53$,
$-0.09$).  For non-diagonal families, gradient optimisation narrows the
margin gap without closing it: Jordan near-boundary improves from
$-1.86$ (scalar) to $-0.97$; nonnormal stable from $-11.07$ to
$-4.98$.

Regulariser comparison (Experiment~2) shows all methods achieve 100\%
post-hoc certification in the standard constrained regime.  The spectral
penalty achieves the lowest task loss ($\leq 6\%$ overhead), the Lyapunov
penalty incurs up to 24\% overhead but achieves the highest resolvent margin
($+0.92$), and the contour barrier provides the best spectral-radius
preservation ($\varrho \approx 0.68$) by not penalising stable modes beyond
the certificate threshold.

Adversarial regime (Experiment~3) reveals the fundamental limitation of
hinge regularisers: at lr = 0.1 with unconstrained radii, 90\% of
unregularised runs escape the unit disk, and hinge regularisers reduce
this to at best 80\% failure (Lyapunov, 20\% stable) because a single
Adam step overshoots the unit circle before the relu gradient fires.
Constrained parameterisation ($\texttt{clamp}$ on log-radius) eliminates
instability entirely by construction.

%% ============================================================
\section{Discussion}
\label{sec:discussion}

\paragraph{Reference optimisation.}
Gradient optimisation improves mean resolvent margins by $2.2\times$
over eig-shrunk on diagonal matrices ($+0.96$ vs.\ $+0.43$) and raises
certification from 36\% to 100\% on trained SSMs vs.\ random-search.
For non-diagonal structures the diagonal reference class is insufficient
regardless of optimisation method; structured non-diagonal references
(e.g.\ Schur-form based) are a natural extension.

\paragraph{Contour barrier vs spectral/Lyapunov.}
All three regularisers achieve 100\% certification on the diagonal SSM
tasks tested, matching the unregularised baseline.  The spectral penalty
achieves the lowest task-loss overhead ($\leq 6\%$) and the highest
resolvent margin among easy-to-implement methods.  The contour barrier's
distinctive advantage is $\varrho$ preservation: by firing only when the
resolvent margin falls below $\tau$, it avoids penalising stable modes
that are well within the certificate, preserving $\varrho \approx 0.68$
vs.\ $0.46$--$0.50$ for spectral/Lyapunov.  This matters for long-memory
AR tasks where near-boundary modes carry useful predictive signal.

\paragraph{Hinge regularisers under adversarial gradient pressure.}
Experiment~3 establishes that relu-based hinges cannot prevent
instability when learning rates are high enough that a single gradient
step overshoots the unit circle.  Softplus hinges would provide
non-zero gradient even far below threshold, but at the cost of a
constant background penalty that accumulates across all modes.  The
fundamental bottleneck is not the hinge form but the step-size: for
Adam with lr$=0.1$, momentum can carry eigenvalues $\Delta\varrho \approx
0.15$ in one step, far exceeding the distance from the near-boundary
init ($\varrho=0.96$) to instability.  Gradient clipping or a smaller
lr would mitigate this; constrained parameterisation eliminates it.

\paragraph{Limitations.}
The gradient-optimised reference is computed on the \emph{initial}
model and frozen; co-optimising $A_0$ jointly with $A$ during training
is a natural extension deferred to future work.
The diagonal-fast-path assumption ($A$ diagonal) is inherent to the
current diagonal complex SSM architecture; extension to structured
non-diagonal SSMs requires the full $O(Kn^2)$ SVD path.
The DLR extension (Section~\ref{sec:dlr}) converges more slowly than
the diagonal path due to the non-convex stability penalty and additional
parameters; a tighter constrained reparameterisation for the low-rank
correction is an open problem.

%% ============================================================
\section{Conclusion}
\label{sec:conclusion}

Gradient-based reference optimisation substantially closes the
certification gap identified in Paper~1.  In the standard constrained
regime, all three regularisers achieve 100\% post-hoc certification;
they differ only in the task-loss--spectral-radius trade-off.  Under
adversarial gradient pressure (unconstrained radii, high lr), hinge
regularisers reduce but cannot eliminate instability: only constrained
parameterisation provides a hard stability guarantee.  The DLR reference
extension demonstrates that richer reference classes are reachable by
gradient optimisation, motivating future work on tight non-diagonal
reparameterisations.

%% ============================================================
\bibliographystyle{plain}
\bibliography{refs}

\end{document}
